\section{The original half-level and an integral PEL model}
\label{sec:integral-pel-frontier}

The cyclotomic field already prescribed by the source resolves the
prime-to-residue-characteristic neatness issue.  It also lets us apply
ramified PEL compactification theorems to the original lattice and
primitive line.  This supplies the source's individual \texttt{Model}
object; its stronger simultaneous model, section and uniformization
contracts are stated separately below.

Retain the host, datum and notation of
Section~\ref{sec:generic-pel-frontier}.  Put $L=\operatorname{lcm}(M,A)$,
and write $K=c.F$ for an original coefficient field.  In particular
$K$ already contains the prescribed root
$\zeta_L=e^{2\pi i/L}$.  No coefficient-field extension is made.
Write $\mathcal O=\mathcal O_E$ and $\overline{\mathfrak b}$ for the
conjugate of the original half-level ideal.

\begin{theorem}[Completing the original finite level]
\label{thm:original-half-level-completion}
Over $K$, the original half-$M$, full-$A$ moduli problem is equivalent
to the moduli problem with the full, symplectically liftable $L$-level
marking and the prescribed Weil-pairing root $\zeta_L$.
The equivalence preserves the abelian scheme, polarization, coefficient
action, rigidified cubical line, and both original label maps.
Forgetting the added labels is an equivalence of functors, not a
cover of positive degree.
\end{theorem}
\begin{proof}
The coefficient idempotents modulo $M$ give complementary submodules
\[
 V_+=\overline{\mathfrak b}\Lambda/M\Lambda,
 \qquad V_-=\mathfrak b\Lambda/M\Lambda,
 \qquad \Lambda/M\Lambda=V_+\oplus V_-.
\]
Both are isotropic for $B$ modulo $M$.  For example, conjugate
linearity of the adjoint moves one factor of
$\overline{\mathfrak b}$ to $\mathfrak b$, so a pairing of two vectors
in $\overline{\mathfrak b}\Lambda$ is divisible by $M$.
Since $B$ is perfect, the cross-pairing $V_+\times V_-$ is perfect.
The same decomposition and perfect cross-pairing hold in
$\mathcal A[M]$ by the original geometric full-Tate condition.

The given half-level restricts to an isomorphism on $V_+$.
Cartier duality and the fixed root $\zeta_M=\zeta_L^{L/M}$ therefore
give a unique marking on $V_-$ with the prescribed cross-pairing.
These are morphisms of finite \'etale group schemes, so the construction
is functorial over arbitrary $K$-schemes.  It gives a unique full
$M$-marking $\alpha_M$.

Let $d=\gcd(M,A)$.  The markings $\alpha_M$ and the given
$\alpha_A$ agree on $d$-division.  This can be checked on geometric
points: an original coherent Tate marking has root
$\zeta'_A=\zeta_A$, hence
$\zeta'_d=\zeta_A^{A/d}=\zeta_M^{M/d}$.  It agrees with the
half-marking, and duality uniquely determines the missing factor
modulo $d$.  The two compatible markings consequently glue to a
unique full $L$-marking.  Explicitly, the map
\[
 \mathcal A[L]\longrightarrow
 \mathcal A[M]\times\mathcal A[A],\qquad
 x\longmapsto([L/M]x,[L/A]x)
\]
identifies its source with the pairs having the same image in
$\mathcal A[d]$.  This assertion is checked prime by prime on the
constant torsion modules.  It also proves uniqueness and agreement
with both entire original label maps.

We must still verify full-Tate liftability; a finite symplectic
marking alone would not suffice.  At a geometric point choose the
original coherent Tate marking $\alpha$ and roots $\zeta'_n$.
There is an $r\in\widehat{\mathbb Z}^{\times}$ with
\[
  (\zeta'_L)^r=\zeta_L,\qquad r\equiv1\pmod A.
\]
We claim there is a unit $u\in\widehat{\mathcal O}^{\times}$ such
that $u\bar u=r$, $u\equiv1\pmod A$, and $u$ acts as the identity
on $V_+$.  At every finite place of $F$ above a rational prime
dividing $M$, the extension $E/F$ is split: otherwise its unique
prime would be fixed by conjugation, contradicting
$\mathfrak b+\overline{\mathfrak b}=\mathcal O$ and
$\mathfrak b\overline{\mathfrak b}=(M)$.
On each split pair use $(1,r)$, with the first component corresponding
to $V_+$.  This also respects the full $A$-level, even when $M$ and
$A$ share prime powers.

At the other places use surjectivity of the norm on units in an
unramified quadratic extension.  If $p^a\mid A$, use its
surjectivity from $1+p^a\mathcal O_{E,v}$ onto
$1+p^a\mathcal O_{F,v}$.  To see the latter directly, write
$p^a\mathcal O_{F,v}=(\varpi^e)$ and lift successively: on every
quotient $1+\varpi^j\mathcal O$ modulo
$1+\varpi^{j+1}\mathcal O$, the norm is the residue-field trace,
which is onto.  Completeness gives a compatible unit.  The unrestricted
unit assertion starts with the finite-field norm and then uses the
same lifting.  This includes residue characteristic two and ramified
$F/\mathbb Q$.  At split places the assertion is immediate.

Now $\alpha\circ u$ has the chosen $L$-root, preserves the original
half-$M$ and full-$A$ labels, and is a full Tate marking.  Its finite
$L$-marking must be the unique one constructed above.  Thus the
original liftability condition is preserved.  Restriction gives the
inverse functor, and all constructions commute with base change and
isomorphisms.
\end{proof}

\begin{corollary}[The exact neat subgroup]
\label{cor:original-neat-level}
Let $H$ be the integral similitude stabilizer of the original finite
labels, with multiplier $\nu$.  Then
\[
 H_0:=H\cap\{g:\nu(g)\equiv1\pmod M\}
       =U(L),
\]
the full principal congruence subgroup of level $L$.  If the original
\texttt{GoodAt} condition holds at $p$, namely
$M/p^{v_p(M)}\geq3$, then the prime-to-$p$ image $H_0^p$ is neat.
\end{corollary}
\begin{proof}
For $g\in H_0$, fixing the half-labels gives
$g-1,g^{-1}-1\in\mathfrak b\operatorname{End}_{\widehat{\mathcal O}}
(\widehat\Lambda)$.  The integral adjoint satisfies
$g^*=\nu(g)g^{-1}$.  Since $\nu(g)-1\in M\widehat{\mathbb Z}$,
we have $g^*-1\in\mathfrak b\operatorname{End}(\widehat\Lambda)$.
Taking adjoints gives
$g-1\in\overline{\mathfrak b}\operatorname{End}(\widehat\Lambda)$.
Intersection gives $g\equiv1\pmod M$, and the given full $A$-level
gives $g\equiv1\pmod L$.  Conversely $U(L)\subset H_0$.
The prime-to-$p$ image is contained in principal level
$M/p^{v_p(M)}\geq3$, so Serre's congruence criterion applies
\cite[Remark 1.4.1.9]{integral-pel-lan}.
This uses the cyclotomic component equivalence just proved; actual
automorphism rigidity by itself would not imply neatness of $H$.
\end{proof}

\begin{theorem}[The prescribed individual integral model]
\label{thm:original-integral-pel-model}
Let $v$ be a finite place of the original coefficient field $K$, with
residue characteristic $p$, and suppose $\mathrm{GoodAt}(p)$.
There is a normal flat projective scheme $\mathcal X/\mathcal O_{K,v}$
containing the original $X_K$ as a dense open subscheme of its generic
fibre, with the original semiabelian $\mathcal O_E$-family, rigidified
symmetric cubical line and primitive identity.  Its generic-fibre
abelian locus is exactly $X_K$.  Thus it supplies every field in the
source's individual \texttt{Model} contract over $\mathcal O_{K,v}$.
The marking comparison and the line comparison both have their
original rigidifications; no power of the line or rescaling of its
generic section lattice is introduced.
\end{theorem}
\begin{proof}
If the datum has rank zero, use $\operatorname{Spec}\mathcal O_{K,v}$
and the trivial group and line.  Otherwise use the original perfect
alternating lattice and PEL datum from
Theorem~\ref{thm:original-generic-pel-moduli}.  Its pairing convention
is $-2\pi iB$, and the fixed norm-minus-one adelic unit from that
theorem recovers the original labels exactly.  As that unit is
coefficient-scalar, it commutes with $U(L)$; the level and neatness
computed above are unchanged.

Choose a projective compatible rational polyhedral cone decomposition
$\Sigma$, which exists by
\cite[Proposition 2.8]{integral-pel-projective}.  Use the singleton
lattice collection $(1,\Lambda,B)$, without an isogeny twist.
Lan's ramified compactification theorem
\cite[Theorem 6.1]{integral-pel-projective} applies to $U(L)$ since
its prime-to-$p$ image is neat.  It gives a normal flat projective
scheme over the localization of the reflex ring at $p$, with its
tautological semiabelian family and coefficient action.  The reflex
field is contained in the source's $\mathbb Q(\sqrt{-d})\subset K$.

Base change to $\mathcal O_{K,v}$, take the reduced closure of the
generic cyclotomic fibre with root $\zeta_L$ (retaining all its
connected components), and normalize.
The normalization is finite because the schemes are of finite type
over an excellent discrete valuation ring.  The result is projective
and normal.  Each of its components dominates the base; hence it is
torsion free, and therefore flat over the discrete valuation ring.
Its generic fibre is the original full-level toroidal compactification
on that component, because this generic fibre is already normal.
Theorem~\ref{thm:original-half-level-completion} identifies its interior
with precisely $X_K$.  If a \texttt{Moduli} object was already chosen,
its original reduced-base representing property uniquely identifies
it with this smooth moduli scheme: apply that property to the two
reduced universal bases and compose the resulting maps.
Pull back the tautological semiabelian family and its coefficient
action; smoothness and the semiabelian fibre condition are preserved
under this base change, including normalization.

It remains to extend the actual line.  The primitive line
$\mathcal L_w$ initially lives on the dense interior $X_K$, rather
than on the entire generic toroidal fibre, and is symmetric and
cubical there.  For a semiabelian scheme
over a normal base, restriction of cubical lines to any dense open is
fully faithful, and every symmetric generic cubical line extends
\cite[Proposition 3.3.2.2 and Theorem 3.3.2.3(2)]{integral-pel-lan}.
Thus $\mathcal L_w$ has a unique extension as a cubical line.
The cubical structure includes its zero rigidification.  Symmetry,
the coefficient adjoint identities and
\[
 \mathcal L_w\simeq
 (\operatorname{id},\iota(w))^*\delta\mathcal L_w
\]
extend by full faithfulness.  Their relations hold because they do
on the dense generic fibre.  This is exactly the original
\texttt{SAB} line package, not merely its polarization or its square.

Finally the formal boundary stratum indexed by a nonzero cusp has a
positive-rank torus with the prescribed character lattice, whereas
the zero cusp is the abelian interior
\cite[Theorem 6.1(3)--(4)]{integral-pel-projective}
\cite[Section 6.2.5]{integral-pel-lan}.
Consequently the generic abelian locus is precisely the original
open moduli scheme.  These assertions supply, respectively, the
source's normality, projectivity, flatness, dense generic open,
semiabelian line data, exact generic identification and
\texttt{ab\_gen} fields.
\end{proof}

This closes an actual arithmetic model construction subobligation.
It does not yet prove the stronger \texttt{LevelStructureConditions},
their simultaneous finite level transitions and global domination,
the rank-one integral section condition \texttt{PushInvertible}, or
every field of \texttt{PELUnif}/\texttt{TorsionIso}.  Lan's theorem
does give formal identifications with the original-lattice Mumford
families, including ramified primes
\cite[Theorem 11.2]{integral-pel-ramified}; checking the full coefficient,
period and section reconstruction interfaces is the next geometric
frontier.  The global theta denominator, section remainder, cusp
comparison and global norm estimates therefore remain open.
